\section{Fundamentos linguísticos e revisão dos conceitos}
\subsection{Mudança diacrônica não se reduz à derivação}
Mudança linguística designa transformações em sistemas e práticas de linguagem ao longo do tempo. Pode afetar pronúncia, contrastes, estruturas morfológicas, construções sintáticas, significados, usos pragmáticos e convenções de escrita. ``Evolução'' significa aqui mudança histórica, sem supor progresso, aperfeiçoamento ou seleção biológica literal. A diversidade dos mecanismos é central nos estudos históricos \cite{campbell2021,martelotta2008}.

Na morfologia sincrônica, \emph{raiz} e \emph{radical} dizem respeito à estrutura de unidades analisadas em determinado sistema, com definições dependentes da teoria. A \emph{etimologia} investiga a origem histórica. \emph{Derivação morfológica} é um mecanismo de formação de palavras, como a combinação de base e afixo; não é requisito para que uma palavra tenha história. A descendência de um item pode envolver erosão e reanálise que tornem a estrutura ancestral opaca. A capacidade de receber sufixos não é teste necessário nem suficiente de continuidade.

Correspondências sonoras regulares ajudam a reconstruir relações; analogia, empréstimo e mudança semântica também precisam ser considerados. Não é necessário observar diretamente cada geração intermediária. A reconstrução se apoia na convergência de evidências, não numa gravação contínua de todos os estágios. A metáfora do ``DNA dos vocábulos'' tem limites: palavras não possuem material genético biológico e podem atravessar fronteiras entre línguas e suportes.

\subsection{Fonética, fonologia, grafia e nome de letra}
A fonética examina propriedades articulatórias, acústicas e perceptivas; a fonologia investiga a organização dos contrastes sonoros. A grafia tem convenções próprias e não é transcrição acústica transparente. O sinal \forma{k} pode ativar o nome da letra, normalmente \forma{cá} em português, mas não é tecnicamente correto chamar \forma{cá} de ``fonema da letra'': é uma sílaba composta por segmentos. Letra, som e fonema precisam ser distinguidos \cite{seara2015}.

A consoante /k/ é uma oclusiva velar surda, articulada com fechamento pelo dorso da língua na região do véu palatino. A vogal acompanhante não integra esse fonema. As letras \forma{c}, diante de \forma{a}, e o dígrafo \forma{qu}, em muitos contextos, podem representar /k/. Disso não resulta identidade fonética integral entre \forma{cá}, \forma{quá} e \forma{quiá}. Nem toda gargalhada contém oclusivas velares. A classificação de /b/ como dental ou alveolar, presente no material inicial, é incorreta: /b/ é bilabial.

Síncope é perda no interior de uma forma; apócope, perda no final; aférese, perda no início. A supressão inicial em \forma{você} $>$ \forma{cê} é aférese, não apócope da primeira sílaba. Metátese, assimilação, redução e palatalização são processos distintos. A ``lei do menor esforço'' pode ser uma intuição sobre tendências articulatórias, mas não explica sozinha a propagação social de inovações.

\subsection{Onomatopeia, iconicidade e convencionalização}
Onomatopeia é representação linguisticamente organizada de som. Seu caráter icônico é uma relação percebida de semelhança entre forma e conteúdo. A semelhança é seletiva e culturalmente mediada: sistemas fonológicos, repertórios gráficos e práticas de leitura participam da interpretação. Iconicidade e convenção não são incompatíveis. Iida e Akita \cite{iida2024} mostram que familiaridade e experiência na língua influenciam julgamentos de iconicidade.

Uma expressão imitativa pode ser criada num episódio e depois repetida por outros usuários. Pode ser emprestada, adaptada ou incorporada a processos morfológicos. Portanto, ``é onomatopeia'' não implica ``não tem etimologia''. Ocorrências interjetivas podem apresentar autonomia prosódica ou sintática relativa, o que não significa ausência de organização linguística. Uma convenção digital pode ser histórica mesmo quando sua transformação é principalmente gráfica.

\emph{Coincidência fonética} é uma hipótese sobre a relação entre formas: elas se parecem sem que a semelhança decorra de descendência ou empréstimo demonstrável. Pode haver coincidência casual ou convergência favorecida pelo mesmo estímulo acústico. Não se devem fundir coincidência e onomatopeia numa categoria. Uma palavra não onomatopaica pode coincidir com outra; duas onomatopeias podem ter relação histórica.

\begin{table}[htbp]\centering\small
\caption{Dimensões que substituem a oposição rígida entre ``evolução real'' e ``onomatopeia''.}
\begin{tabularx}{\textwidth}{p{0.19\textwidth}YY}\toprule
Dimensão & Questão pertinente & Inferência indevida\\\midrule
História & Há transmissão ou empréstimo documentável? & Semelhança prova descendência.\\
Motivação & Há semelhança percebida com um som? & Iconicidade exclui convenção.\\
Morfologia & A forma participa de construções ou famílias? & Sem afixação não há história.\\
Pragmática & Que ação realiza na interação? & Todo marcador implica riso físico.\\
Suporte & Que recursos e restrições são relevantes? & O teclado determina a inovação.\\
Documentação & Qual é a fonte, a data e o localizador? & Não encontrar prova ausência.\\\bottomrule
\end{tabularx}\end{table}

\subsection{Continuidade lexical em três idiomas}
O português \forma{olho} exemplifica continuidade lexical a partir do latim \forma{oculum}, convencionalmente citado no acusativo. Sua história envolve redução e transformação consonantal, com palatalização no percurso românico. Não se reduz à ``sonorização de oc-''. \forma{Voar}, relacionado a \forma{volare}, e \forma{lua}, relacionado a \forma{luna}, ilustram perdas intervocálicas em contextos pertinentes, sem uma receita irrestrita para todas as palavras \cite{ilari2018,campbell2021}.

No inglês, \forma{good} remonta ao inglês antigo \forma{g\=od}. O exemplo combina continuidade e mudança de pronúncia, sem grande alteração gráfica. Sua trajetória não deve ser atribuída mecanicamente à Grande Mudança Vocálica: histórias particulares incluem encurtamentos e condicionamentos. No francês, \forma{eau} é relacionado ao latim \forma{aqua}, com forte divergência gráfica e sonora; a etimologia é registrada pela Académie française \cite{gooddict,academieeau}. Nenhum exemplo exige reconhecimento consciente da origem pelos falantes atuais.

\begin{table}[htbp]\centering\small
\caption{Três exemplos de continuidade lexical.}
\begin{tabularx}{\textwidth}{p{0.14\textwidth}p{0.24\textwidth}Y}\toprule
Idioma & Relação histórica & Cautela\\\midrule
Português & \forma{oculum} $>$ \forma{olho} & Fundamenta-se na história românica, não só no mesmo referente.\\
Inglês & \forma{g\=od} $>$ \forma{good} & Mudanças particulares não se reduzem a uma única lei geral.\\
Francês & \forma{aqua} $>$ \forma{eau} & A relação pode persistir com pouca semelhança superficial.\\\bottomrule
\end{tabularx}\end{table}

\subsection{\forma{Você}: gramaticalização e história social}
\forma{Vossa mercê} $>$ \forma{você} é redução e reorganização histórica de uma expressão de tratamento, associada à pronominalização. A expressão-fonte contém possessivo e nome; a forma resultante desempenha funções de referência ao interlocutor. A análise envolve categoria, distribuição e significado social. Não é um exemplo direto de palavra latina que apenas preservou seu radical: o ponto de partida imediato é uma construção portuguesa.

A cronologia ``vossa mercê no XVI, vosmecê no XIX, você no XX e cê no XXI'' é inadequada. Monte \cite{monte2019} documenta \forma{você} em 1638, refutando seu suposto nascimento no século XX. As variantes não se sucedem como degraus exclusivos: coexistem, distribuem-se regionalmente e adquirem especializações pragmáticas. Eventos como a chegada da corte em 1808 podem integrar hipóteses sociais, mas não substituem documentos datados.

A concordância de \forma{você vai} conserva propriedades da construção nominal de tratamento em sua integração ao sistema pronominal. A referência é ao interlocutor, enquanto a concordância padrão tem forma de terceira pessoa. Não é uma ``anomalia''. A distribuição de \forma{tu} e \forma{você} também varia no Brasil. Sua história articula estrutura linguística, polidez, circulação social e organização colonial \cite{monte2019}.

\subsection{\forma{Vilão}: semântica, estrutura e acepções coexistentes}
\forma{Vilão}, historicamente relacionado a \forma{villanus} e à base \forma{villa}, é pertinente à pejoração: denominações sociais podem adquirir avaliações negativas. Entretanto, uma narrativa única de ``camponês livre'', ``pessoa rude'' e ``antagonista'' não descreve todas as épocas e regiões. Estatutos sociais e acepções dependem do contexto; sentidos antigos e recentes podem coexistir.

A explicação formal \forma{-us} $>$ \forma{-ão} é incorreta. A história das terminações e da nasalidade precisa considerar a base em \forma{-anu-} e a evolução românica. No plano gramatical, \forma{vilão} pode funcionar como substantivo e adjetivo: um sentido moral mais abstrato não obriga a mudança de classe. Categoria, significado e forma não mudam necessariamente em bloco \cite{ilari2018,campbell2021}.

\subsection{Comparações internacionais e falsas etimologias}
O latido em \forma{au-au}, \forma{woof} ou \forma{bow-wow} e o choro em \forma{buá} ou \forma{ouin} mostram diversidade de convenções, mas não demonstram sozinhos independência histórica. Famílias linguísticas diferentes não impedem empréstimos, inclusive de formas icônicas. É preciso investigar cada item.

Um contraexemplo ao material inicial é o japonês \forma{chikutaku}: o verbete reunido em Kotobank identifica relação com o inglês \forma{ticktack} \cite{kotobank}. Apresentá-lo como prova de criação isolada entre Japão e Europa é impróprio. Sem análise especializada, não se adotam como fatos uma origem imitativa do mandarim \forma{gou}, um suposto inglês antigo \forma{gaw}, uma leitura egípcia \forma{iu-iu} ou uma origem onomatopaica consensual do latim \forma{basium}. Etimologia remota incerta não prova imitação independente.

Formas como o tailandês \forma{555} e a repetição coreana de uma consoante ilustram recursos gráficos diferentes: nome de numeral e sinal consonantal. A comparação tipológica não exige afirmar isolamento cultural. A transliteração ``chau-chau'' atribuída à forma coreana no material inicial não é empregada como descrição fonética.

As menções a Ælfric, Chaucer, Shakespeare e Aristófanes permanecem pistas bibliográficas, não testemunhos autenticados neste corpus. O coaxar em \textit{As rãs} não deve ser reclassificado como riso; a grafia grega contemporânea não pode ser projetada sobre a Grécia antiga sem edição e análise. A alegação do ``primeiro haha do Ocidente'' exige investigação especializada e não é adotada como resultado.
